% =======================================================================================
\section{Experiments on XS}
\label{app:xs}
% =======================================================================================

\subsection{Setup}

XS is the smallest rung: $R=16$, 64-bin windows with a 48-bin interior, $f_0$ between 0.1 and
1.26\,mHz (53\% of the log band), four sources per window and 32 data tokens
(\cref{tab:rungs}). Because a step is cheap (39\,ms at width 512), XS is where every design
decision was made. With a tight compute budget, we never ran single-variable sweeps. Each
run combined the changes we believed in, read against the runs already paid for, and was
followed by cheap evaluation-only tests that ruled out alternative explanations. The five
models of \cref{tab:xs-ladder} form the ladder of this section:
\begin{enumerate}
  \item \textbf{uniform clock}: width 512, uniform clock ($p=1$), auxiliary heads over the
        first half of training ($f_{\rm aux}=0.5$), constant learning rate, 500k steps then
        continued to 1M;
  \item \textbf{warp}: the warped clock ($p=3$) and a short auxiliary phase
        ($f_{\rm aux}=0.1$), otherwise identical, 1M steps;
  \item \textbf{warp, 768}: the same at width 768 (170M parameters);
  \item \textbf{warp, cooldown}: model 2 continued for 200k steps with the learning rate
        decaying linearly to zero;
  \item \textbf{warp, 768, cooldown} (\emph{final}): model 3 continued the same way.
\end{enumerate}
The evaluation is that of \cref{app:evaluation}: 840 sources in 210 windows, scored
identically for every model.

\begin{table}[h]
\caption{The XS ladder, on the same 840 sources. Loud width: median $f_0$ width of the found
sources at SNR $\ge100$, and its median ratio to the ideal width $\sqrt3/(\pi\snr)$. Found:
fraction of sources with width below one bin. Calibration: coverage of the 68\% and 95\%
intervals and mean rank of the truth at SNR $\ge100$. Bias: the largest median offset across
eight slices of the source's position in the window, SNR $\ge40$ (\cref{fig:bias}). Hours: A100
time, cooldowns added to the run they continue.}
\label{tab:xs-ladder}
\centering
\small
\setlength{\tabcolsep}{4pt}
% generated by paper/scripts/tables.py from the scorecards; do not edit by hand
\begin{tabular}{llllcccccccccl}
\toprule
& & & & \multicolumn{2}{c}{loud $f_0$ width ($\rho\geq100$)} & \multicolumn{3}{c}{found fraction, $\rho$:} & \multicolumn{3}{c}{calibration, $\rho\geq 100$} & bias & A100 \\
\cmidrule(lr){5-6} \cmidrule(lr){7-9} \cmidrule(lr){10-12}
clock & width & steps & lr & bins & $\times$ ideal & $<15$ & 15--40 & 40--100 & in68 & in95 & rank & [bins] & hours \\
\midrule
uniform & 512 & 1M & const. & 0.147 & 94 & 0.52 & 0.85 & 0.95 & 0.77 & 0.97 & 0.62 & 0.146 & 10.8 \\
warp & 512 & 1M & const. & 0.088 & 59 & 0.45 & 0.81 & 0.92 & 0.87 & 0.99 & 0.50 & 0.045 & 10.9 \\
warp & 768 & 1M & const. & 0.087 & 55 & 0.47 & 0.83 & 0.94 & 0.82 & 0.98 & 0.41 & 0.081 & 16.6 \\
warp & 512 & 1.2M & cooled & 0.039 & 26 & 0.50 & 0.85 & 0.95 & 0.93 & 0.99 & 0.47 & 0.009 & 10.9 + 2.2 \\
warp & 768 & 1.2M & cooled & 0.037 & 25 & 0.54 & 0.86 & 0.97 & 0.93 & 0.99 & 0.49 & 0.008 & 16.6 + 3.6 \\
\bottomrule
\end{tabular}

\end{table}

\subsection{Headline Results}

The final model finds 54\%, 86\%, 97\% and 99.5\% of the sources in the four SNR bands
(\cref{tab:xs-final}), places them without bias (median offsets below 0.006 bins, mean ranks
0.49--0.53), and is calibrated on the conservative side: 78--93\% of the truths inside the
68\% intervals and 97--100\% inside the 95\% intervals. Below SNR 40 its $f_0$ posteriors are
within a factor 2--3 of the ideal width. Above that, the width stops following $1/\snr$
(\cref{fig:width-snr}): it levels off at $0.035$--$0.039$ bins from SNR 100 to 1400, 25 times
the ideal width at the median loud source. This floor, 0.6\,nHz, is the main remaining
limitation, and the rest of this section traces what moved it.

\begin{table}[h]
\caption{The final model (XS, width 768, warped clock, cooldown) by SNR band, on 840 sources.
Width and offset are medians over found sources, in bins of $1/\Tobs$; coverage and rank
over all sources.}
\label{tab:xs-final}
\centering
\small
% generated by paper/scripts/tables.py from the scorecards; do not edit by hand
\begin{tabular}{lrcccrccc}
\toprule
SNR $\rho$ & $n$ & found & width [bins] & $\times$ ideal & offset [bins] & in68 & in95 & rank \\
\midrule
$<15$ & 148 & 0.54 & 0.113 & 2 & -0.006 & 0.78 & 0.97 & 0.53 \\
15--40 & 157 & 0.86 & 0.068 & 3 & -0.003 & 0.90 & 1.00 & 0.50 \\
40--100 & 151 & 0.97 & 0.041 & 5 & +0.000 & 0.90 & 0.99 & 0.49 \\
$\geq 100$ & 384 & 0.99 & 0.037 & 25 & +0.002 & 0.93 & 0.99 & 0.49 \\
\midrule
all & 840 & 0.89 & & & & 0.89 & 0.99 & 0.50 \\
\bottomrule
\end{tabular}

\end{table}

\begin{figure}[h]
  \centering
  \includegraphics[width=0.6\textwidth]{figures/width_snr.pdf}
  \caption{Width of the $f_0$ posterior against the source's SNR, on the 840 scored sources.
  Points: the final model. Lines: medians of the found sources in SNR bins, for the
  uniform-clock baseline, the warped clock, and the final model. Dashed: the ideal width
  $\sqrt3/(\pi\snr)$. Every model has a floor at high SNR; the warped clock lowers it by a
  factor 0.6 and the cooldown by a further 0.44, and below SNR $\sim40$ the final model is
  within a factor 2--3 of the ideal.}
  \label{fig:width-snr}
\end{figure}

\subsection{What Each Change Bought}
\label{app:xs-ladder}

\Cref{fig:ladder} summarises the ladder. Detection differences are paired, source by source,
with a two-sided sign test.

\begin{figure}[h]
  \centering
  \includegraphics[width=\textwidth]{figures/ladder.pdf}
  \caption{The XS ladder. (a) The loud-source floor: median $f_0$ width at SNR $\ge100$, with
  its ratio to the ideal width. (b) Found fraction per SNR band. Same 840 sources for all
  models.}
  \label{fig:ladder}
\end{figure}

\paragraph{Training longer at a constant learning rate helps a little.}
Doubling the uniform-clock run from 500k to 1M steps narrowed the posteriors of the ten eval
windows by a median factor 0.78 (narrower for 27 of 29 sources) and localised two more of the
13 sources at SNR 15--40, but the loud floor only moved from 92 to 76 times the Fisher width.

\paragraph{The warped clock narrows every posterior, at a small cost in detection.}
Against the uniform clock at the same 1M steps, the warp narrowed the loud floor from 0.147 to
0.088 bins, by a median factor of 0.61 at SNR $\ge100$, 0.62 at 40--100 and 0.70 at
15--40, and the posterior was narrower for 97\% of the loud sources. It also removed a bias:
the uniform-clock model placed loud sources low by about 0.056 bins (mean rank 0.62, 8 standard
errors from 1/2). The price was detection: 27 sources were found only by the uniform clock and
6 only by the warp ($p = 3\times10^{-4}$), mostly near SNR 17, the expected effect of moving
training samples from the first half of the path (50\% of them) to its end (21\% left).

\paragraph{Width buys detection, not the floor.}
At a constant learning rate, the 768-wide network had the same floor as the 512 (0.087 against
0.088 bins, flat in SNR), with 2.25 times the parameters and a 1.3\% lower loss. It found 16
sources the 512 missed against 5 the other way ($p=0.027$), recovering half of the warp's
detection cost.

\paragraph{The cooldown halves the floor and removes the bias.}
Cooling the 768 network down over 200k more steps lowered the loss from 0.4499 to 0.4276; at
the constant-rate slope of its last 100k steps the same gain would have taken about a million
more steps. The loud floor fell from 0.087 to 0.037 bins, a median per-source ratio of 0.440
(95\% interval 0.427--0.453), at every loud SNR: 0.038, 0.039, 0.035 and 0.036 bins at SNR
100--200, 200--400, 400--800 and $\ge800$. The cooled model found 21 sources the 1M model did
not, and missed one ($p=10^{-5}$). Against the uniform clock it now finds more sources in every
band (17 against 7). The 512 network gained exactly as much from the same cooldown (factor
0.443), so the gain belongs to the schedule, not to the width. Once both are cooled, the floor
does not depend on width (the 512 is 1.03 times the 768, 95\% interval 1.01--1.05), but the 768
is 11--15\% sharper below SNR 100 and finds 18 sources the 512 misses against 7 ($p=0.043$).

\paragraph{What the floor is not.}
Cheap evaluation-only tests ruled out the obvious numerical causes (\cref{tab:numerics}): the
integrator (32, 64 and 128 RK4 steps agree), bfloat16 at evaluation (float32 on a CPU gives a
width ratio of 0.991), the rounding of the input coordinate at constant learning rate
(\cref{app:geometry-remarks}), and A100 numerics (the A100 and CPU tables agree). Width is ruled
out above. The floor that remains is flat in SNR at 25 times the ideal width. By
\cref{eq:sensitivity} it costs about $2\times10^{-3}$ of the logged loss under the warp, four
times the epoch-to-epoch scatter: visible to the objective, unlike the floor under the uniform
clock, but small next to the 0.022 the cooldown gained. Bfloat16 arithmetic inside the network
during training, and the optimisation itself, are the candidates left untested.

\subsection{Bias}

At a constant learning rate, each run carried its own smooth $f_0$ bias as a function of where
the source sits in the window (\cref{fig:bias}): up to $-0.146$ bins for the uniform clock and
$+0.081$ bins for the 768-wide warp, comparable to their widths, with a different shape in each
run. The cooldown removes it, to at most 0.008 bins across the window, and the mean
offset of the loud sources goes from $+0.029\pm0.005$ to $+0.007\pm0.005$ bins. A profile that
differs from run to run and disappears when the learning rate anneals is the signature of
where a noisy constant-rate iterate happened to stop, rather than of the physics or the data.

\begin{figure}[h]
  \centering
  \includegraphics[width=0.55\textwidth]{figures/bias.pdf}
  \caption{Median $f_0$ offset of the found sources at SNR $\ge40$ in eight slices of their
  position in the window, from its low edge ($-1$) to its high edge ($+1$); bands show one
  standard error of the median. At a constant learning rate each run has its own bias
  profile; after the cooldown none is left.}
  \label{fig:bias}
\end{figure}

\subsection{Calibration}

\Cref{fig:calibration} shows the calibration of the $f_0$ marginal over all 840 sources. The
uniform-clock model's ranks lean to one side, its bias. The final model's are centred but
over-represented in the middle: the truth falls near the centre of the posterior more often
than it should, the signature of posteriors that are wider than necessary, consistent with
in68 = 0.93 at SNR $\ge100$. No model is over-confident in any band. Calibration of the other
parameters will be measured against MCMC (\cref{app:mcmc}).

\begin{figure}[h]
  \centering
  \includegraphics[width=0.55\textwidth]{figures/calibration.pdf}
  \caption{Calibration of the $f_0$ marginal, 840 sources. (a) Empirical CDF of the truth's
  rank among the draws minus the uniform CDF, with a pointwise 95\% band. A calibrated
  posterior stays in the band; the uniform-clock model is biased, and the final model's
  S-shape means its posteriors are too wide. (b) Coverage of the central 68\% (circles) and
  95\% (squares) intervals by SNR band; dashed lines are the nominal levels.}
  \label{fig:calibration}
\end{figure}

\subsection{Single-Source Posteriors}

\Cref{fig:corner-median,fig:corner-loudest} show all eight parameters of two sources,
the median and the loudest of the 40 sources in the eval windows, against the Fisher forecast.
The flow reproduces the exact symmetries: four modes in $\psi$ at quarter-turn spacing, two in
$\phi_0$, and the diagonal bands they form together (\cref{eq:psi-phi-symmetry}). For the
loud source it follows the curved amplitude--inclination degeneracy towards face-on, where
$\psi$ and $\phi_0$ merge into one combination. In both cases the flow is wider than the Fisher
forecast in $f_0$ (the floor above) and along the amplitude--inclination direction. How much
of the latter is the true posterior, which a Gaussian forecast cannot follow, and how much is
excess width, is what the comparison with MCMC will measure (\cref{app:mcmc}). The chirp mass
at these frequencies is not measured, and the flow returns its prior.

\begin{figure}[p]
  \centering
  \includegraphics[width=0.92\textwidth]{figures/corner_median.pdf}
  \caption{Posterior of the median-SNR eval source (SNR 62.4, window at 0.888\,mHz), final
  model, 1024 draws relabelled onto the source (blue), against the Fisher forecast (grey) and
  the truth (black). Contours enclose 39\% and 86\% of the mass. The flow reproduces the four
  $(\psi,\phi_0)$ copies of \cref{eq:psi-phi-symmetry}. 18\% of the draws fall outside the
  plotted $f_0$ range: they are slots the relabelling assigns to this source although they
  belong to an unlocalised SNR-27 source in the same window (\cref{app:evaluation-matching}).}
  \label{fig:corner-median}
\end{figure}

\begin{figure}[p]
  \centering
  \includegraphics[width=0.92\textwidth]{figures/corner_loudest.pdf}
  \caption{As \cref{fig:corner-median}, for the loudest eval source (SNR 1366, window at
  0.101\,mHz). The $f_0$ posterior is far wider than the forecast (the floor); the amplitude
  and inclination follow the degeneracy towards face-on, where the polarisation angle and the
  phase become degenerate.}
  \label{fig:corner-loudest}
\end{figure}

\subsection{Training Dynamics}
\label{app:xs-dynamics}

\Cref{fig:training} shows the training curves. With the original long auxiliary phase, the
flow loss crept down at $2$--$3\times10^{-4}$ per epoch for 170 epochs and then fell by 0.056
within 35 epochs, as the auxiliary weight crossed about 0.2. With $f_{\rm aux}=0.1$ and the
uniform clock (a run stopped after 60 epochs to save compute), the crossover at epoch 36
brought no such drop, so the auxiliary weight alone does not trigger it. Under the warp the
drop came early, near epoch 60, consistent with the warp putting more weight on the fine
structure that the drop seems to learn; what triggers it remains open. After it, the loss
follows a slow power law: each doubling of the steps lowers it by 0.014 to 0.022, and a fit of
$L_\infty + A\,n^{-\alpha}$ to steps 300k--681k predicted 0.4070 at 1M steps, against the 0.4061
reached. The cooldown is worth much more than further constant-rate training:
$-0.022$ in 200k steps for the 768 network and $-0.020$ for the 512. The losses of the
uniform and warped clocks weight the path differently and are not comparable with each other.
Once the auxiliary weight reaches zero, the auxiliary heads stop training and their losses
drift up by orders of magnitude (the reconstruction loss reaches 190 and 520 at 1M steps, at
widths 512 and 768); they play no role in sampling.

\begin{figure}[h]
  \centering
  \includegraphics[width=\textwidth]{figures/training.pdf}
  \caption{Training curves: median flow loss per epoch of 1000 steps. (a) XS, uniform clock,
  500k steps then continued to 1M. (b) XS, warped clock, widths 512 and 768, each continued
  with a 200k-step cooldown (shaded). (c) B, warped clock, 512 wide: the run in progress
  (\cref{app:b}), with its planned cooldown. Losses under different clocks and problems are
  not comparable with each other.}
  \label{fig:training}
\end{figure}

\subsection{Numerical Checks}

\begin{table}[h]
\caption{The warped-clock model (512, 1M steps) under different integrators and precisions.
The float32 rows score the first 200 sources (50 windows) with the same keys; ``ratio'' is
the median per-source width ratio to the default bfloat16 evaluation on the same sources.
Autotuning on the A100 at level 3 gives over-confident posteriors (\cref{app:training-numerics}).}
\label{tab:numerics}
\centering
\small
\setlength{\tabcolsep}{4pt}
% generated by paper/scripts/tables.py from the scorecards; do not edit by hand
\begin{tabular}{lrccccccccc}
\toprule
& & \multicolumn{4}{c}{found fraction, $\rho$:} & \multicolumn{2}{c}{$f_0$ width, $\rho\geq100$} & \multicolumn{2}{c}{$\rho\geq100$} \\
\cmidrule(lr){3-6} \cmidrule(lr){7-8} \cmidrule(lr){9-10}
sampler & $n$ & $<15$ & 15--40 & 40--100 & $\geq100$ & bins & ratio & in68 & in95 \\
\midrule
bf16, A100, 32 RK4 steps (default) & 840 & 0.45 & 0.81 & 0.92 & 0.99 & 0.088 & 1.000 & 0.87 & 0.99 \\
bf16, A100, 64 steps & 840 & 0.45 & 0.81 & 0.92 & 0.99 & 0.089 & 1.000 & 0.86 & 0.99 \\
bf16, A100, 128 steps & 840 & 0.45 & 0.80 & 0.92 & 0.99 & 0.089 & 0.998 & 0.87 & 0.99 \\
fp32, CPU, 32 steps & 200 & 0.39 & 0.78 & 0.94 & 1.00 & 0.088 & 0.991 & 0.85 & 1.00 \\
fp32, A100, autotune level 0 & 200 & 0.39 & 0.78 & 0.94 & 1.00 & 0.088 & 0.991 & 0.85 & 1.00 \\
fp32, A100, autotune level 3 & 200 & 0.45 & 0.83 & 0.97 & 1.00 & 0.036 & 0.411 & 0.45 & 0.88 \\
\bottomrule
\end{tabular}

\end{table}

\subsection{Summary: the Recipe}

What the XS ladder settled, and what B uses: the warped clock with $p=3$, auxiliary heads for
the first 10\% of training, and a linear learning-rate cooldown over the last 20\%. Width 768
helps detection at 1.5 times the cost of 512 at equal steps. The remaining $f_0$ floor sits at
25 times the ideal width for loud sources, with posteriors that are conservative rather than
over-confident.
